\begin{tikzpicture}[
  >={Latex[length=3mm,width=2.2mm]},
  every node/.style={font=\small},
  evec/.style={->,line width=1.3pt,ered},
  pvec/.style={->,line width=1.3pt,vpurple}
]

% Asse verticale con p
\draw[thin,black] (0,-0.5) -- (0,2.5);
\draw[pvec] (0,-0.2) -- (0,1.3) node[left] {$\mathbf{p}$};

% Linea r che va in P
\coordinate (P) at (3.5,2.2);
\draw[thin,black] (0,0.6) -- (P);
\node[midway,above] at (1.6,1.4) {$r$};

% Angolo theta
\draw[thin] (0,0.6) ++(0.6,0) arc (0:33:0.6);
\node at (0.8,0.85) {$\theta$};

% Area giallo chiaro (parallelogramma E1,E2)
\fill[ytangent,draw=ytangent!70!black,opacity=0.8] (P) -- ++(1.8,0) -- ++(0.6,1.4) -- ++(-1.8,0) -- cycle;

% Vettori a P: E orizzontale, E1 verso il basso, E2 diagonale
\draw[evec] (P) -- ++(1.8,0) node[right] {$\mathbf{E}$};
\draw[evec] (P) -- ++(0,-1.4) node[below] {$\mathbf{E}_1$};
\draw[evec] (P) -- ++(1.8,1.4) node[above right] {$\mathbf{E}_2$};

\node[above] at (P) {$P$};

\end{tikzpicture}
